\section{Conclusion}

This thesis studies how systematically weakening excluded-instrument relevance affects the finite-sample performance and informativeness of weak-identification-robust inference in high-dimensional IVQR. The Monte Carlo extension varies a single relevance parameter, $\kappa$, while preserving the marginal treatment distribution, the endogeneity channel, the structural outcome equation, the true quantile treatment effects, the controls, and the observed instrument construction of the baseline design in \textcite{chen_huang_tien_2021}. Oracle-GMM, Full-GMM, and DML-IVQR-BC are evaluated on identical simulated samples across five quantiles and two sample sizes. This design isolates the consequences of progressively reducing the contribution of the excluded instruments to treatment without interpreting $\kappa$ as a complete measure of quantile-specific IVQR identification.

The main result is clear. As excluded-instrument relevance declines, point estimation becomes substantially less accurate and candidate-specific inference becomes less informative. MAE and RMSE increase for all three estimators, including Oracle-GMM. At the median with $n=1000$, for example, Oracle RMSE rises from $0.072$ at $\kappa=1$ to $0.945$ at $\kappa=0.10$, while DML-IVQR-BC rises from $0.104$ to $0.991$. The deterioration of the oracle benchmark shows that this pattern is not generated solely by high-dimensional nuisance estimation. Reduced instrument information directly weakens the ability of the IVQR moments to distinguish the true treatment effect from competing candidate values.

The inference results reinforce this conclusion. The grid-based accepted-set measure within the prespecified domain $A_0=[-1,3]$ expands sharply as relevance falls, while rejection probabilities against false treatment-effect values decline. At the weakest relevance setting, accepted values frequently occupy almost the entire primary domain and power against moderate false alternatives is low even at $n=1000$. Larger samples improve localization and rejection power when usable instrument signal remains, but they do not eliminate the loss of informativeness at $\kappa=0.10$. The magnitude of this deterioration also varies across quantiles: tail quantiles are often more difficult than the median, but there is no universal ordering across all estimators and relevance levels.

A separate finding concerns finite-sample calibration. Empirical coverage is not uniformly close to the nominal $0.95$ level, including in some cells with the strongest relevance setting. Undercoverage also appears for Oracle-GMM, which rules out high-dimensional regularization as the sole explanation. With $R=500$ replications, the maximum Monte Carlo standard error of a reported coverage probability is approximately $0.022$, so the large coverage deviations documented in the simulations cannot be explained by Monte Carlo noise alone. These findings indicate that the asymptotic chi-square reference distribution can be imperfectly calibrated at the simulated sample sizes. At the same time, very high coverage under weak relevance is not necessarily desirable: it can coexist with broad accepted sets and very low power against false candidates.

The estimator comparison clarifies both the usefulness and the limits of DML-IVQR. Regularized nuisance estimation often improves performance relative to unregularized Full-GMM, but DML-IVQR-BC neither dominates in every design cell nor uniformly reproduces the Oracle benchmark. Because the estimators also differ in their residualization implementation, the observed gaps cannot be interpreted as the isolated causal effect of regularization. More fundamentally, Neyman orthogonality protects the target moment against first-order nuisance-estimation error; it does not supply the identifying information lost when excluded instruments become weakly relevant.

The scope of these conclusions is deliberately limited. The analysis uses one preserved DGP, two sample sizes, one same-sample DML-IVQR implementation, a finite numerical candidate domain, and the authors’ asymptotic chi-square inference procedure. It does not study alternative instrument structures, cross-fitted implementations, alternative critical-value procedures, or an empirical application. These are natural directions for further work rather than unresolved requirements of the present experiment.

The central implication is both practical and methodological. In high-dimensional IVQR, successful nuisance regularization and weak-identification-robust test construction do not by themselves guarantee precise or informative inference. When excluded instruments contain little information about treatment, estimation error can rise sharply, false structural effects may be difficult to reject, and finite-sample calibration can remain imperfect. Applied work should therefore evaluate instrument relevance, point-estimation accuracy, coverage, and inferential informativeness separately. No single measure provides sufficient evidence of reliable structural inference.
